% Generated draft: requires mathematical and source review.
\documentclass[11pt]{article}
\usepackage[margin=25mm]{geometry}
\usepackage{fontspec}
\setmainfont{Latin Modern Roman}
\usepackage{amsmath,amssymb,mathtools,unicode-math}
\setmathfont{Latin Modern Math}
\usepackage{enumitem,needspace,xurl,hyperref}
\setlength{\parindent}{0pt}
\setlength{\parskip}{5pt}
\setlength{\emergencystretch}{3em}
\newcommand{\N}{\mathbb{N}}
\newcommand{\Z}{\mathbb{Z}}
\newcommand{\Q}{\mathbb{Q}}
\newcommand{\R}{\mathbb{R}}
\newcommand{\C}{\mathbb{C}}
\newcommand{\F}{\mathbb{F}}
\newcommand{\A}{\mathbb{A}}
\newcommand{\PP}{\mathbb{P}}
\newcommand{\E}{\mathbb{E}}
\newcommand{\eps}{\varepsilon}
\newcommand{\dd}{\,\mathrm{d}}
\begin{document}
\section*{UnsolvedMath ID: 30001522}
\textbf{Source problem code:} catalogue entry
\subsection{Definitions and mathematical statement}
Let \(X\) be a finite connected complex with base point \(x\). Require that \(\pi_1(X,x)\) be abelian and that its natural action on every higher homotopy group be trivial:
\[
\gamma\cdot\alpha=\alpha
\qquad\text{for every }n\geq 2,\ \gamma\in\pi_1(X,x),\ \alpha\in\pi_n(X,x).
\]
The finite-dimensional rational homotopy hypothesis is expressed as
\[
\dim_{\Q}\left(\bigoplus_{n\geq 2}\pi_n(X,x)\otimes_{\Z}\Q\right)<\infty.
\]
Under the usual finite CW complex interpretation, the abelian group \(\pi_1(X,x)\) is finitely generated, so including its rationalization in this direct sum does not change the finiteness condition.

Retain the source notation \(\operatorname{rk}_p(X)\) for free \(p\)-rank and \(\operatorname{rk}_T(X)\) for free toral rank. The supplied record does not define these invariants. In particular, it does not specify whether actions must occur on \(X\) itself or on a space in its homotopy type, which category of actions is allowed, or the precise freeness condition. These conventions must be recovered from the cited report before the statement is fully unambiguous; no definition of either invariant is imposed here. The precise category intended by ``finite complex'' is also unspecified in the supplied record.

Subject to those unresolved source conventions, the question is whether there exists a single such \(X\) for which
\[
\left\{p\in\N : p\text{ is prime and }\operatorname{rk}_p(X)>\operatorname{rk}_T(X)\right\}
\]
is infinite. Thus the same finite complex must satisfy the strict inequality for infinitely many distinct primes.
\subsection{Short English statement}
Does one finite connected complex, with abelian fundamental group acting trivially on all higher homotopy groups and finite-dimensional rational homotopy, have free \(p\)-rank strictly greater than free toral rank for infinitely many primes? The definitions of the two ranks and the intended category of complexes require clarification from the source.
\subsection{Sources}
\begin{itemize}
\item Local UnsolvedMath registry, supplied complete source record: permanent ID \(30001522\), source problem code \texttt{catalogue entry}, title ``Persistent Gaps between Free \(p\)-Rank and Toral Rank.'' The record belongs to the Oberwolfach Reports Open Problems collection, lists proposed year 2010, and supplies no named proposer.
\item Supplied report citation: \emph{Cohomology of Finite Groups: Interactions and Applications} (2011), Oberwolfach Reports, pp.\ 1885--1952, DOI \texttt{10.4171/owr/2010/32}. \url{https://doi.org/10.4171/owr/2010/32}. The supplied literature-source label is ``Oberwolfach Report 32, Cohomology of Finite Groups.''
\item Supplied literature source: Hanke, \emph{The stable free rank of symmetry of products of spheres}. \url{https://doi.org/10.1007/s00222-009-0197-3}.
\item Supplied literature assessment, checked \texttt{2026-08-21}: the record labels the problem open and states that the report explicitly presents this existence question as open and that the cited asymptotic theorem does not settle it. This is supplied dated curation derived from the OWR record, not an independently verified assessment of the literature or current status.
\item The supplied statement status is \texttt{corrected\_verified}; its verification note claims comparison with the official Oberwolfach report and preservation of the hypotheses, notation, and open target. This draft records that supplied claim without independently verifying the report or its mathematical content.
\end{itemize}
\end{document}
